% =====================================================================
%  capitulos/capNN.tex — Plantilla de capítulo
%  Copiar como capNN.tex y agregar \include{capitulos/capNN} en main.tex.
%  Las referencias se filtran por la palabra clave capNN del .bib;
%  las lecturas recomendadas, por lectNN.
% =====================================================================
\capitulo{Título del capítulo}{Nombre del autor}

\section{Introducción}

Texto. Ejemplo numerado del texto (numeración continua en el capítulo):

\begin{exe}
  \ex\label{ex:NN-uno} Oración de ejemplo.
  \ex\label{ex:NN-dos}
    \begin{xlist}
      \ex Variante a.
      \ex[*]{Variante b agramatical.}
    \end{xlist}
\end{exe}

Remisión: \exref{ex:NN-uno}, \exrefs{ex:NN-uno}{ex:NN-dos}.

\begin{ejercicio}
Consigna del ejercicio (la numeración de sus ejemplos se reinicia y va sin
paréntesis).

\begin{exe}
  \ex Primer ítem.
  \ex Segundo ítem.
\end{exe}

\begin{consignas}
  \item Primera pregunta.
  \item Segunda pregunta.
\end{consignas}

\begin{resuelto}
Resolución del ítem 1.
\end{resuelto}
\end{ejercicio}

\begin{recuadro}[Título del recuadro]
Contenido destacado.
\end{recuadro}

% --- Secciones finales -------------------------------------------------
\lecturasrecomendadas

\solucionario
\begin{solucion}{1}
Respuestas.
\end{solucion}

\glosariocapitulo
\begin{glosario}
  \item[Término:] Definición.
\end{glosario}

\referenciascapitulo
